\section{Exam 2.1: Analysis of Decisions}\label{sec:grade10-exam2-1}

\subsection{Assessment and evaluation}\label{sec:grade10-exam2-1-evaluation}

This exam evaluates only the criteria C1--C5. Each criterion is evaluated across the three points below.
A student passes a criterion if it is correctly activated in at least two of the three points.
The overall exam result is binary (Passed/Not passed) and is determined by the set of criteria passed.

\subsection{Point 1}\label{sec:grade10-exam2-1-point1}

A school cafeteria must choose one of three menu plans for a two-week festival: Plan A (local produce focus),
Plan B (mixed menu), or Plan C (bulk staples). Demand for meals can be high, medium, or low.
Based on past festivals, the probability of high demand is $0.35$, medium demand is $0.45$, and low demand is $0.20$.
Profits are measured in hundreds of dollars. If demand is high, the profits are 92 for Plan A, 70 for Plan B, and 55 for Plan C.
If demand is medium, the profits are 50 for Plan A, 62 for Plan B, and 48 for Plan C. If demand is low, the profits are
$-8$ for Plan A, 30 for Plan B, and 38 for Plan C.

\begin{enumerate}[label=\textbf{(C\arabic*)}, leftmargin=2.2em]
	\item Formulate the decision problem by clearly stating the decision alternatives, states of nature, probabilities, and consequences.
	\item Interpret the key elements of the decision problem (alternatives, fortuitous events, and consequences) in this context.
	\item Build the payoff table from the description.
	\item Apply the maximax rule and state the recommended plan.
	\item Apply the maximin rule and state the recommended plan.
\end{enumerate}

\subsection{Point 2}\label{sec:grade10-exam2-1-point2}

A delivery cooperative is choosing a fleet strategy for the next season: Strategy A (electric vans),
Strategy B (hybrid vans), or Strategy C (diesel vans). Demand for deliveries can be strong, steady, or weak.
Fuel costs can be low or high. The cooperative estimates the demand probabilities as $0.30$ (strong), $0.50$ (steady),
$0.20$ (weak). Fuel costs are independent of demand, with probabilities $0.55$ (low) and $0.45$ (high).

Revenue (in thousands of dollars) by demand level is shown below.
\begin{center}
	\begin{tabularx}{0.8\linewidth}{@{}L{0.32\linewidth}C C C@{}}
		\toprule
		Demand level & Strategy A & Strategy B & Strategy C \\
		\midrule
		Strong & 420 & 380 & 340 \\
		Steady & 320 & 300 & 270 \\
		Weak & 210 & 230 & 220 \\
		\bottomrule
	\end{tabularx}
\end{center}

Operating costs (in thousands of dollars) by fuel cost level are shown below.
\begin{center}
	\begin{tabularx}{0.8\linewidth}{@{}L{0.32\linewidth}C C C@{}}
		\toprule
		Fuel cost level & Strategy A & Strategy B & Strategy C \\
		\midrule
		Low & 160 & 190 & 210 \\
		High & 220 & 250 & 290 \\
		\bottomrule
	\end{tabularx}
\end{center}

\begin{enumerate}[label=\textbf{(C\arabic*)}, leftmargin=2.2em]
	\item Formulate the decision problem and define the combined states of nature using the independence assumption.
	\item Interpret the decision alternatives, states of nature, and consequences for this setting.
	\item Build the payoff table by computing profits for each combined state and the probability of each state.
	\item Apply the maximax rule and state the recommended strategy.
	\item Apply the maximin rule and state the recommended strategy.
\end{enumerate}

\subsection{Point 3}\label{sec:grade10-exam2-1-point3}

A community sports center must choose a membership model for the next year: Model A (standard monthly pass),
Model B (premium pass with classes), or Model C (flex pass). Demand for memberships can be high, medium, or low.
Staffing costs can be favorable or unfavorable. Demand probabilities are $0.25$ (high), $0.50$ (medium),
$0.25$ (low). Staffing costs are independent of demand, with probabilities $0.60$ (favorable) and $0.40$ (unfavorable).

Expected number of members by demand level is shown below.
\begin{center}
	\begin{tabularx}{0.85\linewidth}{@{}L{0.32\linewidth}C C C@{}}
		\toprule
		Demand level & Model A & Model B & Model C \\
		\midrule
		High & 620 & 480 & 700 \\
		Medium & 460 & 360 & 520 \\
		Low & 300 & 240 & 360 \\
		\bottomrule
	\end{tabularx}
\end{center}

Membership fees are $35$ (Model A), $48$ (Model B), and $30$ (Model C) dollars per member per year.
Variable staffing cost per member is $18$ when costs are favorable and $26$ when costs are unfavorable.
Fixed annual costs are $6{,}500$ for Model A, $9{,}000$ for Model B, and $5{,}200$ for Model C.
All monetary amounts in this point can be treated in dollars.

\begin{enumerate}[label=\textbf{(C\arabic*)}, leftmargin=2.2em]
	\item Formulate the decision problem, specifying how states of nature are formed and how profits are computed.
	\item Interpret the alternatives, fortuitous events, and consequences in this context.
	\item Build the payoff table by calculating profit for each combined state and assigning the correct probabilities.
	\item Apply the maximax rule and state the recommended model.
	\item Apply the maximin rule and state the recommended model.
\end{enumerate}
